% ============================================================
% Capitolo 6 — Discussione (traduzione italiana di
% chapters/ch06_discussion.tex; label invariati)
% ============================================================
\chapter{Discussione}
\label{ch:discussion}

Questo capitolo interpreta i risultati del \cref{ch:results}: che cosa
l'ordinamento dei task ha cambiato e che cosa no (\cref{sec:ordering}),
che cosa ha contribuito la cornice a mondo condiviso
(\cref{sec:ptc-revisited}), e le minacce e i limiti di scope sotto i
quali ogni affermazione vale (\cref{sec:threats}, \cref{sec:scope}).

% --- §6.1–6.2: interpretazione della matrice completa. ---
\section{Che cosa cambia l'ordinamento dei task (e che cosa no)}
\label{sec:ordering}

Leggendo insieme i quattro assi di confronto del \cref{ch:results},
l'ordinamento dei task ha fatto qualcosa di diverso da ciò che la
letteratura sui curricula di solito promette e da ciò che l'intuizione
sull'addestramento misto di solito teme.

\paragraph{Non la velocità di apprendimento.}
L'asse su cui i curricula sono classicamente attesi ripagare ---
l'efficienza campionaria iniziale --- non ha mostrato nulla: tutte e
quattro le celle hanno superato le stesse soglie di SR a finestra entro
pochi percento dello stesso conteggio di passi
(\cref{sec:sample-efficiency}). Una lettura plausibile è che il design
della reward di questa piattaforma risolva già il problema per cui i
curricula vengono di solito ingaggiati: con bonus densi per waypoint e
shaping sulla distanza (\cref{sec:rewards}), anche i percorsi da
\SI{100}{m} forniscono segnale di apprendimento fin dai primi episodi,
quindi i task facili non servono per l'avvio. Da curricula costruiti
per l'esplorazione a reward sparse non ci si dovrebbe aspettare il loro
vantaggio caratteristico dove la reward è densa --- una condizione di
scope su H1 più che una confutazione del curriculum learning in
generale.

\paragraph{Selezione dell'equilibrio.}
Ciò che l'ordinamento ha cambiato è \emph{quale policy lo stesso
apprendista è diventato}. La meccanica è visibile nei vincoli:
completare un percorso da \SI{100}{m} entro l'orizzonte di 1{.}000
passi richiede una velocità media sana, quindi una policy esposta ai
percorsi hard dal passo~0 deve guidare veloce per riuscire del tutto
--- e paga in collisioni (26\% in valutazione per MLP-batch). Una
policy che spende il suo primo quarto di addestramento su percorsi da
\SI{30}{m} può riuscire lentamente, e cementa la cautela (10.2~km/h,
fallimenti dominati dallo stallo). L'ordinamento dei task, in questa
lettura, agisce da \emph{selettore di equilibri} sul trade-off
velocità--sicurezza più che da acceleratore. Le celle GNN sostengono la
lettura selettiva dall'altro lato: con uno stimatore di valore diverso,
lo stesso curriculum ha selezionato un equilibrio diverso (più veloce)
(\cref{sec:defensive-case}) --- il regime orienta la scelta del bacino,
ma non la determina da solo.

\paragraph{Non uno scudo contro la degradazione.}
L'erosione entro la run della \cref{tab:quartiles} è di piattaforma:
colpisce entrambi i regimi, a livello di task fissato, in entrambe le
classi. Qualunque sia il suo meccanismo --- co-adattazione tra classi e
annealing dell'entropia sono i candidati a registro --- l'ordinamento
non la causa né la previene. L'implicazione pratica per chi riusa
questa piattaforma è che i numeri cumulativi \emph{finali} sottostimano
i picchi di metà addestramento in ogni configurazione, e la selezione
del checkpoint durante l'addestramento cambierebbe le conclusioni ---
una scelta di reporting che questa tesi ha deliberatamente evitato in
favore di endpoint a budget fisso.

\paragraph{Generalizzazione: il beneficio portabile.}
L'unico vantaggio del curriculum replicato in entrambe le architetture
del critic è il margine su Town05 ($+19.0$ e $+10.0$\,pp,
\cref{tab:interaction}) --- ottenuto nonostante la copertura in-domain
più ampia del batch, e contro l'intuizione
copertura-implica-robustezza di H2. Un meccanismo provvisorio coerente
con entrambe le coppie: la policy batch si co-adatta per l'intero
budget a una singola distribuzione stazionaria (una mappa, tutte le
lunghezze), mentre la policy curriculum deve restare competente
attraverso una sequenza \emph{mutevole} di distribuzioni --- una forma
blanda di regolarizzazione implicita per non-stazionarietà. A $n{=}1$
per cella questa resta un'ipotesi; ciò che i dati stabiliscono è la
direzione e la sua replicazione tra architetture.

\paragraph{E i pedoni.}
Per la classe più debole e strutturalmente vincolata
(\cref{sec:ped-limits}), l'ordinamento è stato irrilevante in tre celle
su quattro; l'eccezione (GNN-curriculum, $+12$\,pp in valutazione)
suggerisce che i pedoni beneficino solo quando il critic può
rappresentare la struttura tra agenti \emph{e} l'esposizione è
graduale. Come lezione di design: le scelte di regime calibrate sulla
classe forte non vanno assunte trasferibili alla debole.

% ------------------------------------------------------------
\section{La copertura parallela dei task rivisitata}
\label{sec:ptc-revisited}

La cornice a mondo condiviso (\cref{sec:parallel-coverage}) è stata
adottata per il throughput --- sei istanze di task per episodio --- con
campioni correlati e non-stazionarietà reciproca accettati come costi.
Col senno di poi, la campagna mostra quei ``costi'' fare lavoro
scientifico di prim'ordine.

\paragraph{Le interazioni erano portanti.}
L'equilibrio difensivo --- il fenomeno comportamentale più
consequenziale del registro di sviluppo e della cella MLP-curriculum
--- esiste \emph{solo} perché le reward dei veicoli sono condizionate a
un canale di hazard alimentato dal comportamento dei pedoni
(\cref{sec:rewards}). Il risultato architetturale speculare punta nella
stessa direzione: sostituire un critic piatto con uno che modella gli
agenti come nodi interagenti ha cambiato gli esiti in ogni cella
(\cref{sec:arch-axis}), cosa che non avrebbe potuto fare se la
struttura tra agenti fosse trascurabile. In un design a un agente per
mondo, nulla di tutto ciò sarebbe stato visibile.

\paragraph{La correlazione è reale ed è stata trattata come tale.}
Sei record per episodio condividono uno stato del mondo; trattarli come
indipendenti sovrastimerebbe l'evidenza. La tesi non si appoggia quindi
mai ai conteggi di record per l'inferenza (\cref{sec:threats}), e il
fallimento della pipeline di valutazione della \cref{sec:bugs} --- dove
il collasso dei seed del mondo ha reso tutti gli episodi
quasi identici e prodotto un tasso di successo piatto e troppo pulito
--- è il caso limite ammonitore della stessa correlazione spinta alla
degenerazione.

\paragraph{La non-stazionarietà ha una firma visibile.}
Ogni agente si addestra contro cinque co-giocatori in movimento.
L'erosione entro la run controllata per livello della
\cref{tab:quartiles}, più forte esattamente dove la stima del valore è
più debole (celle MLP) e dimezzata sotto il critic relazionale, è il
pattern che la deriva da co-adattazione produrrebbe --- benché il
design non possa isolarla da altri effetti di fine addestramento
(\cref{sec:threats}).

\paragraph{Che cosa la cornice non può rivendicare.}
Senza un controllo a un agente per mondo (\cref{sec:scope}), il
beneficio di throughput della copertura parallela dei task è un
argomento di design, non una misura; e ogni risultato qui sopra è
condizionato a questo regime ricco di interazioni. La cornice va letta
come l'ambiente \emph{entro il quale} il confronto tra regimi vale, non
come una tecnica validata indipendentemente.

% --- §6.3–6.4: indipendenti dai risultati. ---
\section{Minacce alla validità}
\label{sec:threats}

\paragraph{Non-determinismo del simulatore in condizioni identiche.}
La minaccia dominante, quantificata anziché assunta: tre run a
lunghezza piena con codice, configurazione e seed byte-identici hanno
prodotto tassi di successo cumulativi dei veicoli tra 52.2 e 62.4\%
(\cref{sec:iterations}). CARLA sotto carico multi-agente non è
deterministico da run a run nemmeno in modalità sincrona con tutti i
seed fissati. Ogni effetto in questa tesi è quindi letto contro questa
banda di $\sim$10\,pp: le differenze al suo interno non sono mai
attribuite a un fattore sperimentale, qualunque ne sia la direzione. La
banda è stata misurata sul trunk pre-campagna a livello di
addestramento; la si assume rappresentativa per le run dell'era della
campagna, il che è a sua volta un'approssimazione.

\paragraph{Dimensione campionaria.}
Ogni cella contiene una sola run: il design è stato ridotto da tre seed
appaiati per cella a uno per ragioni di calcolo, una decisione
dichiarata nella \cref{sec:campaign}. Ogni conclusione poggia quindi su
contrasti appaiati, non su distribuzioni stimate. L'appaiamento dei
seed e le configurazioni byte-identiche rendono i contrasti
internamente puliti, e il contrasto primario porta con sé una
replicazione incorporata --- è osservato indipendentemente sotto
entrambe le architetture del critic --- ma nessuna inferenza statistica
formale è possibile a questo $n$, e nessuna è rivendicata. Le
mitigazioni sono la banda di repliche di cui sopra e le verifiche di
coerenza --- accordo di direzione tra livelli di difficoltà, scenari di
valutazione, classi di agenti e strati di misura --- anziché test di
significatività.

\paragraph{Rumore di valutazione.}
La valutazione finale usa 100 episodi per scenario, cioè 300 record
veicolo; a tassi di successo di fascia media il solo errore standard
binomiale è $\approx$3\,pp per scenario, prima di considerare la
correlazione entro episodio dei tre veicoli, che riduce ulteriormente
il campione effettivo. Le differenze a livello di singolo scenario di
pochi punti non sono quindi individualmente significative; si
interpretano solo i pattern tra scenari.

\paragraph{Accoppiamento reward--curriculum.}
La funzione di reward è stata calibrata attraverso iterazioni di gate
eseguite per lo più su configurazioni bloccate su easy
(\cref{sec:gate}), quindi lo shaping del trunk può essere implicitamente
sbilanciato verso i percorsi corti. Entrambi i regimi condividono la
reward identica, quindi il contrasto primario resta internamente
valido; l'avvertenza riguarda l'interpretazione della prestazione
\emph{assoluta} sui percorsi hard, e la possibilità che una reward
calibrata diversamente sposterebbe i due regimi in modo diverso.

\paragraph{Singola coppia di mappe.}
La generalizzazione è misurata su esattamente un trasferimento, Town03
$\to$ Town05. I due layout differiscono in natura
(\cref{sec:carla-bg}), il che rende il test significativo, ma le
affermazioni sono limitate a questa coppia; non si offre evidenza che
l'ordinamento dei regimi si trasferisca ad altre mappe, composizioni di
traffico o condizioni meteo.

\paragraph{Storia della misura.}
La pipeline di valutazione stessa ha fallito la validità una volta
(\cref{sec:bugs}), e il difetto di lunghezza dei percorsi ha
ridefinito silenziosamente la distribuzione dei task per un'intera era
di esperimenti. Entrambi sono stati colti da segnali di integrità
(output implausibilmente regolari; log per agente mancanti). Il
protocollo ora in vigore --- verifiche di integrità, configurazioni
appaiate, ricalcolo dai log grezzi --- è stato costruito in risposta,
ma la lezione generale resta come classe di minaccia: in una
piattaforma custom, gli effetti apparenti devono prima sopravvivere
alla domanda ``lo strumento sta misurando ciò che dichiara?''

\paragraph{Fattibilità dei percorsi pedonali.}
I risultati dei pedoni a difficoltà medium/hard sono confusi da un
limite strutturale dell'ambiente: la topologia dei marciapiedi di
Town03 spesso non può sostenere i target nominali dei percorsi
(\cref{sec:routes}). La \cref{sec:ped-limits} separa il fallimento
attribuibile all'ambiente da quello attribuibile alla policy, e le
conclusioni di regime lato pedoni sono ristrette di conseguenza.

% ------------------------------------------------------------
\section{Limitazioni di scope}
\label{sec:scope}

Oltre alle minacce alla validità interna, il design fissa
deliberatamente diverse dimensioni; i risultati non vanno letti oltre
di esse.

\begin{itemize}
  \item \textbf{L'architettura degli actor è fissa.} Entrambe le classi
        usano actor MLP ovunque; varia solo l'\emph{encoder del critic}
        (MLP vs.\ GNN, solo lato addestramento). Questa tesi non fa
        alcuna affermazione sull'architettura della rete di policy, e
        l'asse GNN parla solo della stima del valore sotto CTDE.
  \item \textbf{La cornice a mondo condiviso non è ablata.} La
        copertura parallela dei task (\cref{sec:parallel-coverage}) è
        la cornice fissa di entrambi i bracci, non un trattamento: non
        c'è un controllo a 1 agente per mondo, quindi il suo contributo
        causale all'efficienza campionaria è caratterizzato, non
        isolato.
  \item \textbf{Osservazioni privilegiate, niente percezione.} Gli
        agenti osservano feature di stato curate
        (\cref{sec:obs-actions}), non sensori. Tutti i risultati
        riguardano la decisione sotto interazione; nulla segue per
        pipeline basate su camera o lidar.
  \item \textbf{Un algoritmo, una versione della piattaforma.} MAPPO su
        RLlib~2.10.0 e CARLA~0.9.16 ovunque. Gli effetti di regime sono
        condizionati a questo stack; altri algoritmi possono interagire
        diversamente con l'ordinamento dei task.
  \item \textbf{La difficoltà è solo lunghezza del percorso.} La
        configurazione definisce famiglie di difficoltà per densità di
        traffico e miste, ma sono inutilizzate (\cref{sec:campaign});
        ``difficoltà'' in ogni affermazione significa lunghezza target
        del percorso a traffico fisso.
  \item \textbf{Popolazione fissa.} Sempre 3 veicoli $+$ 3 pedoni;
        nessuna evidenza di scaling in nessuna direzione.
  \item \textbf{Generazione dei percorsi pedonali.} Il planner a catene
        di marciapiede con il suo tetto agli attraversamenti è parte
        della piattaforma studiata; i suoi limiti noti di fattibilità
        delimitano ciò che le policy pedonali possono esprimere,
        indipendentemente dal regime di addestramento.
\end{itemize}
